\chapter{A risk figure is not a decision}\label{ch:decision}

Most people who work with confidential statistics hold a belief that seems too plain to
need stating: once the disclosure risk of a release has been measured carefully, the
decision follows. If the measured risk is low, the release is safe; if it is high, the
file needs more protection. The hard work, on this view, is the measurement.

The belief has a strong case behind it. The methods of statistical disclosure control
have been developed over four decades, are documented in handbooks, implemented in
tested software and taught to the staff who apply them
\citep{willenborg2001,hundepool2012,dwork2014}. Cell suppression, minimum cell counts,
global recoding, record swapping, added noise, synthetic data and formal privacy
mechanisms are all well understood. An institution that measures risk with these tools,
and releases only when the measured figure is small, appears to have done everything a
careful institution could do. What could a committee add to a well-measured number
except its signature?

Two facts break the belief, and both come from the evidence reported in this book. The
first is that one release has no single risk figure. On the simulated release of
Chapter~\ref{ch:measures}, three measures that are routinely reported under the one
label ``re-identification risk'' give 0.466, 0.064 and 0.023.
% source: paper/generated/numbers.tex (prosecutor mean 0.466, population-frequency version 0.064, unconditional 0.023)
Each is correct, and each describes a different experiment. A committee that reads ``0.023''
has not measured the risk; it has measured the answer to one question, and nothing on the
page says which. The second fact is that even a well-defined figure is an estimate. For a
release whose true risk sits exactly at a tolerance of 0.12, a simulation of 1,200 attacks
produces an estimate at or below the tolerance with probability 0.52.
% computed in Chapter 4 from the Binomial(1200, 0.12) distribution
A committee that releases whenever the estimate is below the line therefore releases a
boundary file about half the time, and its records, which hold only the estimate, cannot
show this.

This chapter asks one question, and the rest of the book answers it in detail: \emph{what
must an institution add to a measured risk before the result can be called a decision?}

\section{A Thursday committee}

Consider the statistics release committee of the Bank, the fictional central bank whose
staff appear throughout this book. On Thursdays the committee reviews the files proposed
for release. Naledi, the disclosure control officer, presents a table for each file: the
share of sample-unique records, the smallest cell count, a mean re-identification
probability. The chair asks whether anyone objects, and if nobody does the file goes out.

The procedure works much of the time, because the people in the room know the data and
the users. Its weakness shows only later. A year after one such meeting, Grace from
internal audit is asked to check a release that a journalist has written about. She finds
the minutes and Naledi's table. She cannot find what adversary the table assumed, what
probability of harm the committee meant to accept, who accepted it, or when the
assessment was meant to lapse. The decision may well have been right. Nobody can show
that it was.

\section{Three objects run together}

Grace's difficulty has a precise cause. A release decision involves three different
objects, and practice treats them as one. The first is a probability model for an
attempted linkage: given what a declared adversary knows and what the release shows, how
likely is it that the adversary names the right record. The second is a statistical
estimate of that probability, which in all but the simplest cases comes from simulating
the attack many times and so carries sampling error. The third is a policy decision: how
much residual risk is acceptable, for which harm, on whose authority and until when. The separation is not new. \citet{duncan1986} treated disclosure as a change in what an intruder can infer, which ties every risk figure to a model of the intruder, and statistical decision theory has long kept the model of the world, the evidence about it and the loss attached to each action apart \citep{berger1985}.

The first is a modelling problem, the second an inference problem and the third a
governance problem. Each needs its own discipline. Running them together produces failures
that are easy to miss, because each looks reasonable from inside one of the three.

The first failure concerns the model. A risk figure means nothing until it states its
adversary: who attacks, with what knowledge, against which targets, and whether the target
is assumed to be in the file. The published cases show how much the adversary's knowledge decides. \citet{sweeney2002} linked hospital records with names removed to a voter list through date of birth, sex and postcode, and \citet{narayanan2008} matched film ratings in a released file to reviews their authors had posted in public. Neither attack would have shown up in a table of key frequencies computed without the outside source. Chapter~\ref{ch:scenario} derives the most familiar measure
in the field, the reciprocal of the key frequency, from four assumptions, and shows that
two of them fail in the ordinary case. A release that keeps ten records in every key class
can still hold a record that an adversary with one extra variable identifies with
probability 0.74.
% source: paper/generated/table_kanonymity.tex; paper/bis_ifc_full.tex (Section 10)

The second failure concerns the estimate, and it is the coin toss described above.
Chapter~\ref{ch:estimate} shows that replacing the point estimate with a one-sided upper
confidence bound controls the chance of a wrong release at a stated level, in the sense of a significance test \citep{neyman1933}, and Chapter~\ref{ch:selection} extends the control to the choice among several protection settings. A small case shows the size of the difference.

\begin{example}[A boundary file on a small budget]\label{c1:ex-boundary}
The numbers are illustrative. Suppose Naledi can afford only 100 simulated attacks on a file whose true risk is exactly the committee's tolerance of $0.12$. A rule that releases when the estimate is at or below $0.12$ releases whenever 12 or fewer attacks succeed, and at a true risk of $0.12$ that happens with probability $0.58$. Suppose the run gives 11 successes. The estimate is $0.11$, below the line. The exact one-sided 95 per cent upper bound of \citet{clopper1934} for 11 successes in 100 attacks is $0.176$, well above it. The bound rule passes this file only if 6 or fewer attacks succeed, since 6 successes give a bound of $0.115$ and 7 give $0.127$, and at a true risk of $0.12$ that happens with probability $0.037$. The first rule releases the boundary file more often than not. The second releases it fewer than four times in a hundred.
\end{example}

The third failure concerns the decision. An institution that says ``we apply a minimum
cell count of ten'' has stated a protection setting and called it a policy. It has not
stated the probability of harm it accepted, so nobody can check whether the setting
achieves the policy, and nobody owns the judgement that the residual risk is acceptable.
When the data or the users change there is no record of what the setting was meant to
deliver, and so nothing to reassess against. Nor does a cell count say anything about the route by which the data reach their users, which the Five Safes framework treats as a control in its own right \citep{desai2016}.

Naledi met the second failure in a meeting that Chapter~\ref{ch:estimate} recounts. Her
estimate for a household survey extract stood just below the tolerance; her upper bound
stood just above it. The chair asked, reasonably, why a second and larger number should
decide. The answer to that question is the first piece of what an institution must add.

\section{What a decision rule needs}

The book's answer is a written decision rule with three parts, one for each object. The
model part declares an adversary: who attacks, what it knows, which targets it attacks and
which conditioning event the risk figure assumes. The inference part replaces the point
estimate with an upper confidence bound at a stated error probability. The governance part
states a tolerance and attaches to it what turns a number into a decision: an owner who
accepts the residual risk, an expiry date, the harm class the tolerance is set for, the
legal basis and the parties it protects, the access route through which the data will be
used, and, for a series of releases, the horizon over which risk is assessed.
Chapter~\ref{ch:thresholds} sets out these elements and who should write each one.

A release is then admissible when the upper bound for every target in a declared target set
lies at or below the tolerance. The rule is simple, and its simplicity is the point. Two
analysts given the same declarations and the same simulation output reach the same
decision, and an auditor can check afterwards that the decision followed from the evidence.
When Grace reads a record written under this rule, every question she could not answer in
the old minutes has a line of its own.

The rule has a price, and the book does not hide it. The evidence needed to certify a
tolerance grows as the tolerance shrinks. If an attack is simulated many times and never
succeeds, the number of clean attempts needed for a 95 per cent upper bound to reach one in
a hundred is 299. For one in a thousand it is 2,995, and for one in a hundred thousand it is
299,572.
% source: paper/generated/table_certification.tex
Observing no successful attack never shows that the risk is zero, and choosing a tolerance
and choosing the evidence standard turn out to be one decision. The tolerance itself is not
arbitrary either. Chapter~\ref{ch:tolerance} shows that a stated weight on harm against the
value of the data implies a tolerance for each release, and that one tolerance for every
release is usually the wrong choice.

\section{Several releases from one register}

Single releases are not how institutions publish. They publish programmes: a quarterly
series, a research file and a public extract, all drawn from one register. An adversary who
holds several of these is not obliged to use them one at a time. When records keep the same
pseudonym across releases, the adversary can intersect the candidate sets that each release
leaves, and the result can identify a target that no single release exposes.

The consequence is stronger than intuition suggests. Chapter~\ref{ch:programmes} proves that
no function of the per-release risks can bound the risk of an adversary who pools the
releases. On a simulated register, four releases that each keep at least three records in
every key class, and that contain no unique record at all, identify the target with
probability 0.458 once pooled, against an upper bound of 0.253 implied by the four marginal
risks.
% source: paper/generated/numbers.tex; paper/generated/table_programme.tex

At the Bank, this is the problem Tomás brings. He is a visiting researcher who wants a
loan-level file from the credit register for six months, a table of lending by sector and a
public series, for one study. Each request, judged alone, would pass Naledi's assessment. The committee's question is whether the three together do, and its old procedure cannot answer it, because the procedure only ever looks at one file.

\begin{example}[Two tables, one loan]\label{c1:ex-twotables}
The numbers are illustrative. Seven loans appear in two of Tom\'as's files under the same pseudonyms, numbered 1 to 7. The first file gives the borrower's sector: loans 1 to 4 share one sector and loans 5 to 7 another. The second gives the region: loans 1, 5 and 6 share one region and loans 2, 3, 4 and 7 another. Every class in each file holds at least three loans, so neither file has a unique record. An adversary who knows the sector and region of the firm behind loan 1 finds four candidates in the first file and three in the second. Read one at a time, the files give it chances of $0.25$ and $0.33$, and even adding the two, as a bound on succeeding with either file, gives only $0.58$. Pooled, the candidate sets are $\{1,2,3,4\}$ and $\{1,5,6\}$, their intersection holds loan 1 alone, and the adversary names it with certainty.
\end{example}

\section{A running example}

Most chapters return to one example: a synthetic credit register built to share the
structure of a supervisory collection such as the euro area's AnaCredit
\citep{israel2017}. It holds 20,000 firms as borrowers, 60 lenders with unequal shares, 18
sectors, 12 regions and four size classes, observed over six monthly vintages. About 27,000
loans are live at any time. Each released loan carries a persistent pseudonym, its sector,
region and size class, and its exposure.
% source: configs/paper.json; paper/generated/numbers.tex (27,167 loans in the first vintage)

The register is synthetic because real loan-level credit data are not released for general
use, and that choice limits what the example can show. Its parameters are chosen by design,
so the protection settings it implies describe the simulation and not an actual institution.
What it can show is how the decision rule responds when the adversary changes. At an
illustrative tolerance of 0.12, an assessment of one release at a time passes the file with
no suppression. An adversary who pools the six vintages and knows in which of them the
target's loan was live requires a minimum cell count of 75. One who also holds a reading of
the loan's exposure in each vintage requires 150, and at that floor the file keeps 20 per
cent of loans but only 4 per cent of exposure, with no loan to a large firm.
% source: paper/generated/table_robustness.tex
Against that adversary a loan-level file is the wrong product for Tomás's study, and
Chapter~\ref{ch:credit} follows him to the secure room where he does his work instead. Pieter,
a credit analyst at one of the reporting lenders, is the hardest adversary of all, because
he knows his own loans almost exactly; no floor up to 300 stops him.

\section{How the book is arranged}

Part~I sets up the problem: this chapter, the definition of risk as the outcome of a declared
attack (Chapter~\ref{ch:scenario}), and the measures that answer different questions
(Chapter~\ref{ch:measures}). Part~II treats the risk figure as an estimate: its uncertainty,
the choice among settings, and certification while an attack is monitored. Part~III turns to
the decision: the elements of a governance threshold and the tolerance a stated harm weight
implies. Part~IV deals with programmes of releases and with formal privacy. Part~V reports the
evidence from simulation, public microdata, the credit register and a pipeline in which a
model proposes attacks. Part~VI is for practice: falsification tests, an operational protocol,
worksheets. Part~VII follows the attack after the release: an audit that runs while the attack runs (Chapter~\ref{ch:learning}), exposure that changes over time (Chapter~\ref{ch:change}), dynamic models and Bayesian inference (Chapter~\ref{ch:dynamics}), how far the decision moves when its inputs move (Chapter~\ref{ch:chaos}), and attacks planned and recorded so that every retry is counted (Chapter~\ref{ch:planner}). Part~VIII closes with the questions that remain open. Appendix~\ref{app:software} maps the software
that produces every number.

Each chapter is built the same way, and the method comes from research on teaching. \citet{muller2008} found that physics students who believed they already understood a topic learned little from a clear explanation of it, and learned much more from a film that first stated their wrong belief and then broke it. Each chapter of this book therefore opens by naming a belief that a careful reader is likely to hold and stating it in its strongest form. It then gives early the first fact that breaks the belief, and asks one question, which the rest of the chapter answers as a connected case. Two threads run through that case: the technical argument, and the people who meet the problem in practice, either in the published record or at the Bank. The main text ends by answering the question plainly, before the \textit{In practice} block and the proofs.

A reader who must decide on a release next month can read this chapter, the \textit{In
practice} blocks at the end of each chapter, and Part~VI. A reader who wants to check the
claims will find each one stated as a proposition with its assumptions and proved in the last
section of its chapter.

\section{The answer in brief}

The question was what an institution must add to a measured risk before the result can be
called a decision. It must add three things. It must name the attack the figure describes, so
that the figure answers a stated question. It must replace the estimate with an upper bound
at a stated error probability, so that a boundary file is not released by chance. And it must
write down a tolerance with an owner, an expiry date, a harm class, a legal basis, an access
route and, for a programme, a horizon, so that someone answers for the residual risk and
someone else can check the decision afterwards. Without these the committee has a number;
with them it has a decision.

\begin{inpractice}
Before the next release decision, write down four things that are usually left implicit: the
adversary the assessment assumes, the tolerance the release must meet, the person who owns
that tolerance, and the date after which the assessment lapses. If any of the four cannot be
written down, the committee is not yet in a position to say whether the release is safe,
whatever the risk table shows.
\end{inpractice}
